% Titrage pH-métrique de l'aspirine : courbe et dérivée (exercice 13, chapitre 3)
\begin{tikzpicture}
\begin{axis}[otMesures,otAxesBout,name=ph,width=76mm,height=56mm,xmin=0,xmax=25,ymin=0,ymax=13,
  xtick={0,5,...,25},ytick={0,2,...,12},xlabel={$V_\mathrm{B}$ (mL)},ylabel={pH}]
\addplot[otNavy,line width=1pt] table[x index=0,y index=1] {figures/donnees/03-aspirine-ph.dat};
\end{axis}
\begin{axis}[otMesures,otAxesBout,at={(ph.east)},anchor=west,xshift=17mm,width=76mm,height=56mm,xmin=0,xmax=25,ymin=0,ymax=24,
  xtick={0,5,...,25},ytick={0,5,...,20},xlabel={$V_\mathrm{B}$ (mL)},ylabel={$\mathrm{dpH}/\mathrm{d}V_\mathrm{B}$ (mL$^{-1}$)}]
\addplot[otNavy,line width=1pt] table[x index=0,y index=2] {figures/donnees/03-aspirine-ph.dat};
\solution{\draw[otGreenDark,line width=0.6pt,dash pattern=on 2pt off 1.5pt] (axis cs:13.8,0) -- (axis cs:13.8,21.2);
\node[pcLegende,anchor=west,text=otGreenDark] at (axis cs:15,16) {$V_\mathrm{E}=13{,}8$ mL};}
\end{axis}
\end{tikzpicture}
